\section{Maxwellian baseline calibration}
\label{sec:baseline}

We calibrate the reduced model of Section~\ref{sec:model} against
two literature etch-rate anchors for \sfsix{} ICP etching of
silicon. The low-power anchor is taken from \citet{mansano1999}, who
report a mean silicon etch rate of
$1.4\,\mu\mathrm{m\,min^{-1}}$ (i.e.\ $1400\,$nm\,min$^{-1}$) at an
ICP antenna power of $P_L = 200$\,W in pure \sfsix{} at
$10\,$mTorr. The high-power anchor is taken from
\citet{panduranga2019}, who report a mean silicon etch rate of
$2.27\,\mu\mathrm{m\,min^{-1}}$ ($2270\,$nm\,min$^{-1}$) at
$P_H = 2000\,$W in pure \sfsix{} at $30\,$mTorr, zero substrate
bias, 50\,sccm flow, in an Oxford PlasmaPro 100 Cobra300 reactor.
The two anchors span one decade in absorbed power and lie
approximately a factor of 1.6 apart in measured etch rate, the
sub-linear power scaling qualitatively consistent with the onset of
feedstock depletion at the high-power end.

We note two chamber-condition differences between the two anchor
measurements that the model absorbs into its fitted parameters
rather than treating as independent degrees of freedom. First, the
\citet{mansano1999} measurement was made with a $200$\,W ICP
antenna power \emph{and} a $150$\,W RF bias applied to the
cathode, while our model assumes zero bias; the additional
ion-energy-driven wafer contribution is absorbed into the fitted
wafer coupling $\varepsilon$. Second, \citet{panduranga2019}
reports their $2.27\,\mu\mathrm{m\,min^{-1}}$ rate at a chamber
pressure of $30\,$mTorr, whereas the \citet{mansano1999}
measurement and the nominal operating condition of our model are
both at $10\,$mTorr. The identifiability theorem of
Section~\ref{sec:identifiability} is insensitive to this pressure
difference because the invariant
$\mathcal{I} \equiv \beta\,\kdiss\,\nel(P_H)$ is preserved when
pressure-dependent factors move $\nel(P_H)$ and $\beta$ in
compensating directions. We calibrate both anchors at the nominal
$10\,$mTorr gas density, treating the $\sim\!3\times$ pressure
difference as a systematic absorbed into the fitted $\varepsilon$
and $\beta$, and we flag in Section~\ref{sec:discussion} that the
absolute numerical value of $\beta$ should not be read as a literal
gas residence time for any specific reactor.

The calibration procedure follows the two-stage protocol of
Section~\ref{sec:model:calibration}. The first stage sets
$\beta = 0$ and fits $\varepsilon$ against the \citet{mansano1999}
anchor by a single analytical division. The second stage
reintroduces $\beta$ and refits $(\varepsilon,\beta)$ jointly so
that both anchors are reproduced exactly. The two-stage procedure
returns fitted values
$\beta_{\mathrm{Max}} = 14.27\,$ms and
$\varepsilon_{\mathrm{Max}} = 7.73\times10^{-30}$\,m\,atom$^{-1}$.
The subscript denotes the use of Maxwell-convolved rate coefficients
in equations \eqref{eq:power-balance}--\eqref{eq:source},
anticipating the alternative rate source introduced in
Section~\ref{sec:boltzmann}. Both anchors are reproduced to machine
precision: $1400.00\,$nm\,min$^{-1}$ and
$2270.00\,$nm\,min$^{-1}$, respectively. Across the full
intermediate power range from 50 to $3000$\,W, the calibrated model
interpolates smoothly between the two anchors with a sub-linear
power dependence and saturation at the high-power end characteristic
of the depletion correction in equation~\eqref{eq:depletion}. The
numerical behavior of the calibration is as expected for a
one-parameter fit against a two-point dataset, and no
ill-conditioning is detected during the root-finding stage.

The fitted parameters reproduce the calibration data exactly, but
the numerical value of the depletion timescale is not consistent
with any simple physical interpretation. The nominal identification
of $\beta$ with the gas residence time of the chamber implies that
a fitted $\beta = 14.27\,$ms should correspond to a residence time
of the same magnitude. For a cylindrical ICP reactor at
$10\,$mTorr with volume $5$--$20\,$L and inflow
$10$--$100\,$sccm, the geometric residence time
$\tau_{\mathrm{res}} = \nSFsix V / \dot N_{\mathrm{in}}$ lies in the
range $[36,\,1438]\,$ms. The fitted $\beta_{\mathrm{Max}}$ is
smaller than the shortest plausible value by a factor of 2.5 and
smaller than the longest by a factor of more than 100. No
reasonable choice of reactor volume and flow rate at the nominal
calibration condition produces a residence time as short as the
fitted value, and the identification of $\beta$ as a gas residence
time in the Maxwell-calibrated model is therefore untenable.

This discrepancy is the point of departure for the remainder of the
paper. Several explanations are in principle available: the
depletion model may be misspecified, the rate coefficients may be
in error, the reduced model may contain a structural multiplicative
bias, or the calibration procedure may be determining a quantity
different from the one we are interpreting. The next section
analyzes the structure of the two-anchor calibration and establishes
which combinations of parameters are actually constrained by the
anchor data.
